\chapter{Pipeline planning and design}
\label{ch:pipeline-planning}

Analysts need notebooks and ad hoc SQL; ML engineers need scheduled training with promotion gates.
We separate stable contracts (configs, artefacts, lineage) from replaceable infrastructure (orchestrator, object store, registry).

\section{Consumer obligations and guiding principles}
\label{sec:consumers-principles}

Non-functionals: recover partial runs, audit lineage, alert on ingestion or deploy failure.
Stack shape: Airflow-class DAGs, S3-compatible storage, MLflow-style experiment tracking.

\section{Stage coverage: ingestion through monitoring}
\label{sec:pipeline-stages-global}

\begin{figure}[htbp]
  \centering
  \includegraphics[width=\linewidth,height=0.55\textheight,keepaspectratio]{elt.pdf}
  \caption{ELT pattern with object storage.}
  \label{fig:etl-flow}
\end{figure}

\begin{figure}[htbp]
  \centering
  \includegraphics[width=\linewidth,height=0.92\textheight,keepaspectratio]{arch_diag.pdf}
  \caption{End-to-end MLOps pipeline: raw Lichess data lands in MinIO; sequential components (configuration plus prior artefacts $\rightarrow$ processing $\rightarrow$ new artefacts) cover ingestion, validation, transformation, training, and evaluation; a decision gateway permits model push only when evaluation guardrails pass.}
  \label{fig:mlops-pipeline}
\end{figure}

\begin{figure}[htbp]
  \centering
  \includegraphics[width=\linewidth,keepaspectratio]{whowhatandhow.pdf}
  \caption{Pipeline stages: what (inputs and outputs), why (rationale), who (accountability), and how (internal steps and tools) for ingestion through monitoring (Exercise~1, Step~2).}
  \label{fig:pipeline-whowhatandhow}
\end{figure}

\Cref{fig:etl-flow} shows Airflow bucket ETL; \cref{fig:mlops-pipeline} shows immutable artefact stages through validation, training, evaluation, and conditional push; \cref{fig:pipeline-whowhatandhow} tabulates what, why, who, and how for each mandated stage.
The sections below elaborate each stage with dedicated diagrams.

\section{Data ingestion}
\label{sec:stage-ingestion}

\begin{figure}[htbp]
  \centering
  \includegraphics[width=\linewidth,height=0.85\textheight,keepaspectratio]{dataIngestion.pdf}
  \caption{Data ingestion: the orchestrator applies ingestion config, downloads Lichess monthly data to a raw-data bucket, validates checksums, decompresses with Zstandard, parses games, and writes processed Parquet; invalid payloads trigger alerts.}
  \label{fig:data-ingestion}
\end{figure}

Ingestion lands bulk \texttt{.pgn.zst} in MinIO with monthly partitions, checksum contracts, and idempotent keys (\cref{fig:data-ingestion}).
Data engineering owns connectors; governance checks licensing.
Stable landing prevents gaps and duplicates that would break dashboards and labels.

\section{Data pre-processing}
\label{sec:stage-preprocessing}

\begin{figure}[htbp]
  \centering
  \includegraphics[width=\linewidth,height=0.92\textheight,keepaspectratio]{dataValidation.pdf}
  \caption{Data validation under orchestration: train/test splits are checked for schema completeness and drift; Great Expectations stores validation reports; failures route to alerts and updated artefact keys, while validated processed data feeds an offline feature store.}
  \label{fig:data-validation}
\end{figure}

Validation enforces schema and drift budgets; transformation encodes features, splits without leakage, and serialises preprocessors beside tensors (\cref{fig:data-validation}).
SMEs tune rules; automation blocks downstream on threshold breaches.
Analysts consume silver tables; training reuses the same transforms at inference.

\section{Model development}
\label{sec:stage-model-dev}

\begin{figure}[htbp]
  \centering
  \includegraphics[width=\linewidth,height=0.85\textheight,keepaspectratio]{modelTrainer.pdf}
  \caption{Model training: load train/test arrays from validated artefacts, fit the model, compute metrics, register the best run in MLflow, and retain a containerised, reproducible training environment.}
  \label{fig:model-trainer}
\end{figure}

Training reads transform bundles and logs git, configs, and metrics in MLflow (\cref{fig:model-trainer}).
Research approves hypotheses; engineers run sweeps with sequestered test data until final evaluation exports.

\section{Model deployment}
\label{sec:stage-deployment}

\begin{figure}[htbp]
  \centering
  \includegraphics[width=\linewidth,height=0.85\textheight,keepaspectratio]{deployment.pdf}
  \caption{Deployment: promote a recent model from MLflow, build a Dockerised FastAPI image, run health checks, and push to the serving host (e.g.\ EC2); improved metrics can trigger CI/CD to refresh the endpoint.}
  \label{fig:deployment}
\end{figure}

Containers promote only after gates; MLOps owns CI/CD with SecOps image review (\cref{fig:deployment}).
Canary and rollback retag registries; contracts keep serving aligned with preprocessing.

\section{Model monitoring}
\label{sec:stage-monitoring}

\begin{figure}[htbp]
  \centering
  \includegraphics[width=\linewidth,height=0.85\textheight,keepaspectratio]{monitoringObserv.pdf}
  \caption{Monitoring and observability: API metrics scraped by Prometheus, dashboards and automated reports, rule-based routing to Slack/PagerDuty/email, and controlled reruns from user or scheduled inputs.}
  \label{fig:monitoring}
\end{figure}

Production tracks feature drift, latency SLOs, and score error as new months arrive (\cref{fig:monitoring}).
MLOps triages alerts; retrains stay lineage-linked to ingestion batches.

\section{Iterative refinement narrative}
\label{sec:iteration-pipeline}

We replaced an early sketch with~\cref{fig:mlops-pipeline} for validation loops and RACI~\cref{tab:pipeline-raci-lite}.

\begin{table}[htbp]
  \centering
  \caption{Pipeline stage accountability matrix (illustrative).}
  \label{tab:pipeline-raci-lite}
  \footnotesize
  \setlength{\tabcolsep}{4pt}
  \renewcommand{\arraystretch}{1.1}
  \begin{tabular}{@{}>{\raggedright\arraybackslash}p{0.31\linewidth} >{\centering\arraybackslash}p{0.115\linewidth}>{\centering\arraybackslash}p{0.115\linewidth}>{\centering\arraybackslash}p{0.115\linewidth}>{\centering\arraybackslash}p{0.115\linewidth}@{}}
    \toprule
    Pipeline stage & R & A & C & I\\
    \midrule
    Data ingestion             & DE & Head of data & Stewardship & Analysts\\
    Data pre-processing        & DE / MLE & MLOps lead & Domain SMEs & Governance\\
    Model development          & MLE & Research lead & SMEs & Analysts\\
    Model deployment           & MLOps & Prod owner & SecOps & Analysts\\
    Model monitoring           & MLOps & Prod owner & MLE \& DE & Stewardship\\
    \bottomrule
  \end{tabular}
\end{table}

\textbf{Legend:}
\textit{R}esponsible performs the hands-on work,
\textit{A}ccountable approves outcomes,
\textit{C}onsulted offers domain expertise,
\textit{I}nformed stays aware of timelines.
